\section{Pre-Silicon Evaluation of the PiM Read}
\label{sec:results}
\subsection{Setup}
An 8$\times$8 slice of the layout netlist of the macro, with its drivers, word-line flops and level
shifters, drain switches, delay line, and TDC flops, is simulated with the foundry PSP models
(typical corner). Cells are linear resistors:
9\,k$\Omega$ for LRS (the IHP LRS2 target) and, for HRS, 200\,k$\Omega$, the top of the
range IHP has characterized (1\,\si{\micro\ampere} at 0.2\,V); 1\,M$\Omega$ is added as an \emph{ideal case}. At the 0.06--0.12\,V read bias the HRS non-linearity of
these devices raises $R$ by about 10\,\% over its 0.2\,V value and the LRS is ohmic~\cite{zanotti2026model}; the IHP
compact model, at the two states it can produce (1.4 and 12\,k$\Omega$), reproduces the resistor
runs' crossing times within 0.06\,ns at the three Table~\ref{tab:pim} settings.
The supplies are confined to
$\Vd,\Vg,\Vdel\in\{0.8,0.9,1.0\}$\,V so that a read never approaches a write condition, and the
voltage across the selected cells and their current are recorded in every run. An access takes
the precharge (below), the settle, the offset, and the READ pulse: about 0.45\,\si{\micro\second}
for 512 result bits at 100\,MHz, mostly precharge, and 8--13\,pJ per 8$\times$8 slice from the supply currents (65--105\,pJ per access
scaled to the macro, controller excluded). Of this, 90\,\% is the 1.2\,V clock energy of the 64 word-line flip-flops, which toggle every
cycle during the 40-cycle precharge that the extracted loading needs at $\Vd=0.8$\,V; the read path itself is below 0.3\,pJ. A clock gate on the word-line flip-flops during
precharge would remove most of this energy; the figures here are without it.

\textit{Capacitance model.} A resistor cell ignores two capacitances that set the discharge. The
first is the line loading: a Quantus RC extraction of the macro (21\,930 MOS and 4096 MEMRES)
gives 47--62\,fF of interconnect on each BL and 21--36\,fF on each SL; the slice adds 70\,fF to
BL (wire plus the source junctions of the 56 rows it does not instantiate) and 30\,fF to SL. The
second is the MEMRES device itself, a metal-insulator-metal stack that neither the IHP compact
model (whose only state is the filament gap) nor the extraction contains: for its design-rule area of 0.36\,\si{\micro\meter\squared}, the model's
6\,nm HfAlO and $\varepsilon_\mathrm{r}=12$--20, a Sentaurus solution of the drawn stack (Metal2
plate, Via2, Metal3) gives $C_\mathrm{MIM}=6.8$--11.0\,fF per cell, 6\,\% above the parallel plate,
which bypasses a selected cell with a time constant of 2\,ns at 200\,k$\Omega$ and 10\,ns at
1\,M$\Omega$, the scale of the differences the read resolves. All results use
$C_\mathrm{MIM}=10$\,fF; the line
loading, which doubles every crossing time, also makes the default 10-cycle precharge leave BL
at 0.5\,V for $\Vd=0.8$\,V, so all runs precharge for 40 cycles (400\,ns).

\subsection{Single-row read}
The Liberty view of the macro assumes $\Vg=4$\,V, where the cross-coupled level shifters take 89--126\,ns
and the RTL controller in AMS never captured a sample; at 0.8--2.0\,V they settle within
3.7\,ns, so the controller was re-sequenced to open the word lines before the read. At
$\Vg=0.8$\,V the access transistor is in weak inversion and limits the cell current to
0.5--1.7\,\si{\micro\ampere} (200 to 9\,k$\Omega$, 0.9\,V on BL), so the 22:1 ratio of the cells
reaches the bit line as a current contrast of only 3--3.6:1. A single row then separates HRS
from LRS by the full eight taps: the 200\,k$\Omega$ cell crosses the buffer threshold 50.7\,ns
after the drain opens and the LRS cell 28.2\,ns ($\Vd=\Vg=\Vdel=0.8$\,V), a window at +20\,ns
lies between them, and the cell sees 0.09\,V and 1.3\,\si{\micro\ampere}.

\subsection{Multi-row read: NOR and majority}
\label{sec:pim}
\begin{figure}[!t]
\centering
\includegraphics[width=\columnwidth]{grid}
\caption{Multi-row read with the extracted line capacitance and $C_\mathrm{MIM}=10$\,fF, 8 rows
open, $k$ cells at 9\,k$\Omega$ and $8-k$ at the HRS value, 40-cycle precharge, settle-drain,
$\Vg=0.8$\,V. (a) BL crossing time vs.\ $k$ at the best points; shaded: the TDC window at the
+10\,ns read offset (first READ\_SAMPLE pulse), $\Vdel=1.0$\,V. (b) Simulated codes at those
points.}
\label{fig:pim}
\end{figure}
\begin{table}[!t]
\centering
\caption{Best operating point per operation within 0.8--1.0\,V (slice, LRS 9\,k$\Omega$)}
\label{tab:pim}
\scriptsize
\setlength{\tabcolsep}{1.6pt}
\begin{tabular}{@{}lccclcc@{}}
\toprule
op. & HRS & $\Vd/\Vg/\Vdel$ & off. & codes $c(k)$ & margin & $V_\mathrm{RRAM}$/$I_\mathrm{cell}$\\
 & ($\Omega$) & (V) & (ns) & $k=0\ldots8$ & (taps) & (V/\si{\micro\ampere})\\
\midrule
read, 1 row & 200\,k & 0.8/0.8/0.8 & 20 & 8, 0 (HRS, LRS) & 8 & 0.09/1.3\\
NOR, 8 rows & 200\,k & 0.9/0.8/1.0 & 10 & 7 1 0 0 0 0 0 0 0 & 6 & 0.06/0.8\\
majority, 8 rows & 200\,k & 1.0/0.8/1.0 & 10 & 8 8 8 6 3 1 0 0 0 & 2 & 0.06/0.8\\
NOR, 8 rows & 1\,M$^\dagger$ & 0.8/0.8/0.9 & 10 & 8 0 0 0 0 0 0 0 0 & 8 & 0.12/1.0\\
majority, 8 rows & 1\,M$^\dagger$ & 1.0/0.8/1.0 & 10 & 8 8 8 8 7 3 0 0 0 & 4 & 0.12/1.0\\
\midrule
\multicolumn{7}{@{}l}{direct transients, $C_\mathrm{MIM}=10$\,fF, extracted line capacitance; $^\dagger$ideal HRS}\\
\bottomrule
\end{tabular}

\end{table}
The PiM read opens $n$ rows of one slice column in a single access with $k$ cells in LRS and
$n-k$ in HRS. At $\Vd\le1$\,V the settle cycle must ground SL (the settle-drain bit of the
controller): with SL floating, charge sharing takes BL below the buffer threshold before READ. The
discharge starts when the word lines open, and the eight-tap window is placed on it by the
signed read offset in 10\,ns steps. The full grid $\Vd,\Vg,\Vdel\in\{0.8,0.9,1.0\}$\,V was simulated for HRS 200\,k$\Omega$
and 1\,M$\Omega$; codes at every window position follow from the simulated crossing times and the
simulated tap tables, with the comparison
offset fitted on direct transients, and the best point of each operation was then simulated
directly (Table~\ref{tab:pim}).

\textit{Where the window lands.} At 200\,k$\Omega$ every multi-row operation is resolved at
$\Vg=0.8$\,V and at no higher word-line voltage: the current-limited
cells stretch the 8-row discharge to 20--27\,ns (Fig.~\ref{fig:pim}(a)) and $\Vd$ decides where
the band falls against the window grid; at $\Vg\ge0.9$\,V the band is shorter than the 10\,ns step and falls between window positions.
That step, one clock period, is the main limitation of the design, since it caps the row count
and the hot-corner margins below. Over the PSP corners and $-40$ to 125\,\si{\celsius} the 8-row band moves from 16 to
60\,ns against the 10\,ns window step; re-tuning $\Vd$, $\Vg$, $\Vdel$, and the offset per
temperature (Table~\ref{tab:retune}) keeps the single-row read at eight taps and recovers the
8-row NOR at $-40$\,\si{\celsius} (five taps) but not at 125\,\si{\celsius} (one tap), where
the band is shorter than the step.
\begin{table}[!t]
\centering
\caption{Margin (taps) vs.\ temperature, typical corner, 200\,k$\Omega$: Table~\ref{tab:pim} setting fixed at
27\,\si{\celsius} / re-tuned ($\Vd$, $\Vg$, $\Vdel$, offset) per temperature}
\label{tab:retune}
\scriptsize
\setlength{\tabcolsep}{4pt}
\begin{tabular}{@{}lccc@{}}
\toprule
 & $-40$\,\si{\celsius} & 27\,\si{\celsius} & 125\,\si{\celsius}\\
\midrule
read, 1 row & 8 / 8 & 8 / 8 & 8 / 8\\
NOR, 8 rows & 0 / 5 & 5 / 5 & 0 / 1\\
majority, 8 rows & 0 / 2 & 2 / 2 & 0 / 0\\
\bottomrule
\end{tabular}
\end{table}

\textit{NOR.} At $\Vd=0.9$\,V the all-HRS crossing (23.9\,ns) lies inside the window and the
$k=1$ crossing (22.9\,ns) before it: codes 7\,1\,0\,0\,0\,0\,0\,0\,0, a NOR by six taps; the
ideal 1\,M$\Omega$ reaches the full scale.

\textit{Majority.} The $k\to k+1$ step scales as $1/k^2$, so $c(4)-c(5)$ needs the window on the
compressed part of the band: $\Vd=1.0$\,V moves it there, codes 8\,8\,8\,6\,3\,1\,0\,0\,0, a
majority by two taps (four at 1\,M$\Omega$). 

\begin{figure}[!t]
\centering
\includegraphics[width=\columnwidth]{rows}
\caption{Rows per access, $\Vg=0.8$\,V. (a) Best NOR and majority margin vs.\ rows open, over
$\Vd$, $\Vdel$ and window position (1--8 rows: full supply grid; hollow, 16--64 rows: at the best
8-row supplies). (b) Crossing-time step behind each decision, $t(0)-t(1)$ for NOR and
$t(n/2)-t(n/2+1)$ for majority; shaded: the tap pitch at $\Vdel=0.8$--1.0\,V. }
\label{fig:rows}
\end{figure}
\textit{Rows per access.} The same model was run with 1 to 8 rows open over the full supply
grid, and with 16 to 64 rows at the best 8-row supplies (Fig.~\ref{fig:rows}). At
200\,k$\Omega$ the NOR keeps all eight taps up to four rows and four to five taps up to eight;
the majority falls from eight taps at two rows to two or three at six to eight. Every open row adds an HRS current that one LRS cell must stand out against, so the
crossing-time step behind each decision falls roughly as $1/n^2$ (Fig.~\ref{fig:rows}(b)), from
17\,ns (NOR) and 5\,ns (majority) at two rows to 1.4 and 0.4\,ns at eight. From 16 rows on, the 200\,k$\Omega$ step is at or below a tap pitch, and single-tap margins
computed there are tap-edge artifacts. With the ideal 1\,M$\Omega$ the NOR step is still 2.0\,ns at 16
rows, but the window, 1.2--2.5\,ns wide and placed in 10\,ns steps, misses it by 0.03\,ns; at
the IHP HRS one access decides over at most eight rows, and
wider sums are split into eight-row groups combined by the host. Only two rows give a distinct code for every count.

\textit{Read disturb and spread.} A read biases the selected cells in the RESET polarity (BL
high, SL grounded), so repeated reads push toward HRS. $\Vg$, not $\Vd$, sets the stress: at the best points a cell sees at most 0.06--0.12\,V and 0.8--1.3\,\si{\micro\ampere}, below
the 0.2\,V read condition of IHP and a sixth of the 8\,\si{\micro\ampere} RESET target. On these
devices $10^5$ pulses of 10\,\si{\micro\second} at 0.1--0.2\,V (SET polarity) leave every
level unchanged and drift starts at 0.3\,V, while retention alone lowers the LRS by 15\,\% in
100\,h at room temperature~\cite{baroni2026drift}:
the refresh interval is set by retention and temperature rather than by the read
count.
The margins above are at the typical corner with fixed resistances. With a 3--3.6:1
current contrast, a Monte Carlo analysis on the simulated crossing times, checked against 160 SPICE
runs with drawn cell resistances, limits the HRS spread for an
$n$-row decision at $10^{-3}$ error to $\sigma_{\ln R}\approx0.46/\sqrt{n}$ with an ideal
threshold, 0.16 at eight rows and 0.33 at two, against 0.3--0.6 measured on the IHP 250\,nm arrays and 0.5--1 on the
Al-doped 130\,nm devices, whose HRS at the same verify has a median of 90--220\,k$\Omega$ but a
conductive edge pinned at 50\,k$\Omega$~\cite{perez2017hfalo,baroni2025aldoping}; the single-row read tolerates 0.46. A fivefold deeper reset, the 1\,M$\Omega$ rows of
Table~\ref{tab:pim}, buys a full-scale 8-row NOR, a four-tap majority, and an eight-row spread
limit of 0.5, inside the measured range. A deeper reset, a tighter verify window, and fewer rows per access are the levers; a higher
word line adds only 30--45\,\% of contrast. The first two are programming-time costs: a deeper reset takes
more incremental-step-pulse-with-verify (ISPVA) pulses per write (about twelve of 12\,\si{\micro\second} at the 5\,\si{\micro\ampere}
target~\cite{perez2017hfalo}) and, with Al doping, leaves a quarter of the cells at the read
floor~\cite{baroni2025aldoping}; a tighter window costs further verify iterations and the
cells that never enter it.

\textit{Limitations.} All results are pre-silicon, from an 8$\times$8 slice with resistor cells
and a single value of $C_\mathrm{MIM}$, so the absolute crossing times carry the uncertainty of
the extraction and of the TCAD capacitance. The 10\,ns window step limits the logic to eight
rows per access at the IHP HRS and loses it at 125\,\si{\celsius}; a fractional-cycle offset is
the first change for a second version. Process corners enter only through the range of the
8-row band, and the read-disturb bound rests on SET-polarity pulse data, so RESET-polarity
disturb and the HRS spread itself remain to be measured on silicon.
